\section{Training Checkpoints}
\label{app:checkpoints}

The constrained groups and the Standard LLMs overlap only partly in perplexity: the Constrained-BabyLM and BabyLM models are small and trained on 100M words, and few Standard LLMs reach their perplexities. The matched-scale coefficients therefore rest partly on extrapolating the perplexity slope. To fill this range with models of standard architecture and training data, we add the intermediate training checkpoints of two runs. These are the ten deduplicated Pythia models (14M to 12B parameters) at steps 128, 256, 512, 2{,}000, 8{,}000, 16{,}000, 32{,}000, 64{,}000 and 128{,}000 (about 0.27B to 268B training tokens), and the 29 checkpoints of the BabyLM GPT-2 baseline (1M to 1B training words, i.e.\ up to ten epochs of the 100M-word corpus). Step 0 is left out, and the final Pythia checkpoint is the released model already in the benchmark. The Pythia checkpoints belong to the Standard LLM group and the BabyLM GPT-2 checkpoints to the BabyLM group.

Every checkpoint is scored exactly as the released models: surprisal RDM against the 16 reading measures, and whole-network embedding RDM (all layers but layer 0) against the eight neural datasets. Perplexity follows the definition of Section~\ref{sec:design-stats} for checkpoints and released models alike: the exponential of the mean last-layer word surprisal over the unit's words on the reading measures, Podcast, MEG-MASC, Broderick and Pereira, and the corpus perplexity over the sentences on Tuckute, Nieuwland and the two Szewczyk studies. CBR-RNN and CBR-RNN-M have no surprisal on the Broderick and Pereira words and are left out of those two units.

A training run contributes up to nine (Pythia) or 29 (BabyLM GPT-2) rows per dataset unit. Treating these as separate models would count one run many times. The regression therefore adds a random intercept for the model within its lineage to Equation~\ref{eq:group}, so that the checkpoints of a run share one model intercept and one lineage intercept. This intercept had zero estimated variance in every fit of Table~\ref{tab:rsa-ckpt}: once lineage is modeled, models of one family do not differ further beyond what their size and perplexity predict, and the estimates equal those of Equation~\ref{eq:group}. Without the lineage and model intercepts, the hundreds of checkpoint rows dominate the perplexity slope and the constraint coefficients shrink towards zero. The 119 checkpoints (90 Pythia, 29 BabyLM GPT-2) add 119 scored models per dataset unit; the Pythia checkpoints share the model intercept of the released Pythia model of the same run.

\begin{table}[t]
\centering
\small
\setlength{\tabcolsep}{5pt}
\caption{Matched-scale regression with the training checkpoints of Pythia (ten sizes, steps 128--128{,}000) and of the BabyLM
GPT-2 baseline (29 checkpoints, 1M--1B words) added to the original models. Behavior: surprisal RDM, 16 reading measures; brain:
whole-network embedding RDM, the eight neural datasets. Perplexity is computed the same way for checkpoints and original models
(Appendix~\ref{app:checkpoints}); the \emph{Original} columns equal the all-datasets rows of Table~\ref{tab:summary}. Random intercepts for the dataset unit, the model lineage and the
model within its lineage (a training run's checkpoints share one model intercept); REML, Satterthwaite $t$-tests.
Coefficient (SE); reference: Standard LLM. $^{*}p<.05$, $^{**}p<.01$, $^{***}p<.001$.}
\label{tab:rsa-ckpt}
\begin{tabular}{lcccc}
\toprule
 & \multicolumn{2}{c}{Behavior} & \multicolumn{2}{c}{Brain} \\
\cmidrule(lr){2-3}\cmidrule(lr){4-5}
 & Original & + checkpoints & Original & + checkpoints \\
\midrule
Memory constraint $\beta_M$ & $+.028^{**}$ (.011) & $+.029^{**}$ (.011) & $-.189^{*}$ (.092) & $-.089$ (.094) \\
Data constraint $\beta_D$ & $+.030^{*}$ (.014) & $+.036^{*}$ (.015) & $-.261^{*}$ (.116) & $-.220$ (.109) \\
Interaction $\beta_{MD}$ & $-.031^{*}$ (.015) & $-.037^{*}$ (.016) & $+.141$ (.136) & $+.159$ (.138) \\
$\log_{10}$ parameters $\beta_S$ & $-.033^{***}$ (.004) & $-.020^{***}$ (.002) & $+.201^{***}$ (.054) & $+.093^{**}$ (.028) \\
$\log_{10}$ perplexity $\beta_P$ & $-.027^{***}$ (.004) & $-.021^{***}$ (.002) & $-.183^{**}$ (.061) & $-.355^{***}$ (.027) \\
\midrule
Checkpoints added & 0 & 119 & 0 & 119 \\
Rows (of which checkpoints) & 1,870 (0) & 3,774 (1,904) & 934 (0) & 1,886 (952) \\
\bottomrule
\end{tabular}
\end{table}


Table~\ref{tab:rsa-ckpt} compares the original models with the original models plus the checkpoints. The \emph{Original} columns reproduce the all-datasets rows of Table~\ref{tab:summary}: the model-within-lineage intercept has zero variance, and both use the same perplexity definition.

Adding the checkpoints leaves the behavioral coefficients almost unchanged. Both constraints are associated with higher behavioral similarity, with a negative interaction, and Constrained-BabyLM models remain above Standard LLMs ($\beta_M+\beta_D+\beta_{MD}=+.028$, $p<.05$). On the brain side, the penalties shrink and are no longer reliable (memory $-.089$, $p=.34$; data $-.220$, $p=.054$; Constrained-BabyLM $-$ LLM $-.151$, $p=.14$). The perplexity slope doubles ($-.355$ against $-.183$) and the size slope halves. Restricting all models to at most 2B, 1B or 500M parameters gives the same picture: Constrained-BabyLM $-$ LLM contrasts of $+.027$ to $+.039$ on behavior (all $p<.05$) and $-.145$ to $-.178$ on the brain (all $p>.05$).

The checkpoints change the brain result through the perplexity slope. Brain similarity rises faster as perplexity falls during training than it differs between models of different perplexity: splitting perplexity into each model's mean over its checkpoints and the checkpoints' deviations from it gives a within-run slope of $-.394$ and a between-model slope of $-.203$. With this split, the data penalty is reliable again ($-.327$, $p=.014$; Constrained-BabyLM $-$ LLM $-.290$, $p=.017$). A single perplexity slope fitted mostly to training runs thus attributes to perplexity part of the difference between finished models. Whether constrained models are less brain-like than Standard LLMs of the same perplexity therefore depends on which perplexity comparison is meant: between finished models, or along a training run.
